\section{Diffusion Probabilistic Models}
\label{sec:bg}
We review diffusion probabilistic models and  their associated differential equations in this section.
\subsection{Forward Process and Diffusion SDEs}\label{sec:diffusion-sde}
Assume that we have a $D$-dimensional random variable $\x_0\in\R^{D}$ with an unknown distribution $q_0(\x_0)$. Diffusion Probabilistic Models (DPMs)~\citep{sohl2015deep,ho2020denoising,song2020score,kingma2021variational} define  a \textit{forward process} $\{\x_t\}_{t\in[0,T]}$ with $T>0$ starting with $\x_0$, such that for any $t\in [0,T]$, the distribution of $\x_t$ conditioned on $\x_0$ satisfies
\begin{equation}
\label{eqn:q0t}
    q_{0t}(\x_t|\x_0)=\N(\x_t | \alpha(t)\x_0, \sigma^2(t)\Iv),
\end{equation}
where $\alpha(t),\sigma(t) \in \R^{+}$ are differentiable functions of $t$ with bounded derivatives, and we denote them as $\alpha_t,\sigma_t$ for simplicity. The choice for $\alpha_t$ and $\sigma_t$ is referred to as the \textit{noise schedule} of a DPM.
Let $q_t(\x_t)$ denote the marginal distribution of $\x_t$, DPMs choose noise schedules to ensure that $q_T(\x_T)\approx \N(\x_T|\vect{0},\tilde\sigma^2\Iv)$ for some $\tilde\sigma>0$, and the \textit{signal-to-noise-ratio} (SNR) $\alpha_t^2 / \sigma_t^2$ is strictly decreasing w.r.t. $t$~\citep{kingma2021variational}.  Moreover, \citet{kingma2021variational} prove that the following stochastic differential equation (SDE) has the same transition distribution $q_{0t}(\x_t|\x_0)$ as in Eq.~\eqref{eqn:q0t} for any $t\in[0,T]$:
\begin{equation}
\label{eqn:forward_sde}
    \dxv_t = f(t)\x_t \dt + g(t) \dwv_t, \quad \x_0\sim q_0(\x_0),
\end{equation}
where $\wv_t\in\R^D$ is the standard Wiener process, and
\begin{equation}
\label{eqn:ft_gt}
    f(t)= \frac{\dd \log \alpha_t}{\dt}, \quad g^2(t)=\frac{\dd\sigma_t^2}{\dt} - 2\frac{\dd\log\alpha_t}{\dt}\sigma_t^2.
\end{equation}
Under some regularity conditions, \citet{song2020score} show that the forward process in Eq.~\eqref{eqn:forward_sde} has an equivalent \textit{reverse process} from time $T$ to $0$, starting with the marginal distribution $q_T(\x_T)$:
\begin{equation}
\label{eqn:reverse_sde_true}
    \dxv_t = [f(t)\x_t - g^2(t) \nabla_{\x}\log q_t(\x_t)] \dt + g(t) \dd \bar{\wv}_t, \quad \x_T\sim q_T(\x_T),
\end{equation}
where $\bar{\wv}_t$ is a standard Wiener process in the reverse time. The only unknown term in Eq.~\eqref{eqn:reverse_sde_true} is the \textit{score function} $\nabla_{\x}\log q_t(\x_t)$ at each time $t$. In practice, DPMs use a neural network $\epsilonv_\theta(\x_t,t)$ parameterized by $\theta$ to estimate the scaled score function: $-\sigma_t\nabla_{\x}\log q_t(\x_t)$. The parameter $\theta$ is optimized by minimizing the following objective~\citep{ho2020denoising, song2020score}:
\begin{align}
    \Lc(\theta;\omega(t))&\coloneqq \frac{1}{2}\int_0^T \omega(t) \E_{q_t(\x_t)}\Big[\|\epsilonv_\theta(\x_t,t)+\sigma_t\nabla_{\x}\log q_t(\x_t)\|_2^2\Big]\dt \nonumber \\
    &= \frac{1}{2}\int_0^T \omega(t) \E_{q_0(\x_0)}\E_{q(\epsilonv)}\Big[\|\epsilonv_\theta(\x_t,t)-\epsilonv\|_2^2\Big]\dt + C, \nonumber
\end{align}
where $\omega(t)$ is a weighting function, $\epsilonv\sim q(\epsilonv)=\N(\epsilonv|\vect{0},\Iv)$, $\x_t=\alpha_t\x_0+\sigma_t\epsilonv$, and $C$ is a constant independent of $\theta$. As $\epsilonv_\theta(\x_t,t)$ can also be regarded as predicting the Gaussian noise added to $\x_t$, it is usually called the \textit{noise prediction model}. 
Since the ground truth of $\epsilonv_\theta(\x_t,t)$ is $-\sigma_t\nabla_{\x}\log q_t(\x_t)$, DPMs replace the score function in Eq.~\eqref{eqn:reverse_sde_true} by $-\epsilonv_\theta(\x_t,t) / \sigma_t$ and define a parameterized reverse process (\textit{diffusion SDE}) from time $T$ to $0$, starting with $\x_T\sim \N(\vect{0},\tilde\sigma^2\Iv)$:
\begin{equation}
\label{eqn:reverse_sde_theta}
    \dxv_t = \left [f(t)\x_t + \frac{g^2(t)}{\sigma_t} \epsilonv_\theta(\x_t,t)\right] \dt + g(t) \dd \bar{\wv}_t, \quad \x_T\sim \N(\vect{0},\tilde\sigma^2\Iv).
\end{equation}
Samples can be generated from DPMs by solving the diffusion SDE in Eq.~(\ref{eqn:reverse_sde_theta}) with numerical solvers, which discretize the SDE from $T$ to $0$. \citet{song2020score} proved that the traditional ancestral sampling method for DPMs~\citep{ho2020denoising} can be viewed as a first-order SDE solver for Eq.~\eqref{eqn:reverse_sde_theta}. However, these first-order methods usually need hundreds of or thousands of function evaluations to converge~\citep{song2020score}, leading to extremely slow sampling speed.


\subsection{Diffusion (Probability Flow) ODEs}
\label{sec:bg_ode}
When discretizing SDEs, the step size is limited by the randomness of the Wiener process~\citep[Chap. 11]{kloeden1992Numerical}. A large step size (small number of steps) often causes non-convergence, especially in high dimensional spaces. 
For faster sampling, one can consider the 
associated \textit{probability flow ODE}~\citep{song2020score}, which has the same marginal distribution at each time $t$ as that of the SDE. 
Specifically, for DPMs, \citet{song2020score} proved that the probability flow ODE of Eq.~\eqref{eqn:reverse_sde_true} is
\begin{equation}
    \frac{\dxv_t}{\dt} = f(t)\x_t - \frac{1}{2}g^2(t)\nabla_{\x}\log q_t(\x_t), \quad \x_T\sim q_T(\x_T),
\end{equation}
where the marginal distribution of $\x_t$ is also $q_t(\x_t)$. By replacing the score function with the noise prediction model, \citet{song2020score} defined the following parameterized ODE (\textit{diffusion ODE}):
\begin{equation}
\label{eqn:diffusion_ode}
    \frac{\dxv_t}{\dt} = \hv_\theta(\x_t,t)\coloneqq f(t)\x_t + \frac{g^2(t)}{2\sigma_t}\epsilonv_\theta(\x_t,t), \quad \x_T\sim \N(\vect{0},\tilde\sigma^2\Iv).
\end{equation}
Samples can be drawn by solving the ODE from $T$ to $0$.
Comparing with SDEs, ODEs can be solved with larger step sizes as they have no randomness. Furthermore, we can take advantage of efficient numerical ODE solvers to accelerate the sampling. 
\citet{song2020score} used the RK45 ODE solver~\citep{dormand1980family}  for the diffusion ODEs, which generates samples in $\sim$ 60 function evaluations to reach comparable quality with a 1000-step SDE solver for Eq.~\eqref{eqn:reverse_sde_theta} on the CIFAR-10 dataset~\citep{Krizhevsky09learningmultiple}. However, existing general-purpose ODE solvers still cannot generate satisfactory samples in the few-step ($\sim$ 10 steps) sampling regime. To the best of our knowledge, there is still a lack of training-free samplers for DPMs in the few-step sampling regime, and the sampling speed of DPMs is still a critical issue. 



